\chapter{OBJETIVOS}

\section{Objetivo Geral}

O presente estudo tem como objetivo principal desenvolver e validar um 
modelo de AM capaz de prever o desempenho de 
equipes de futebol a partir de dados estatísticos de desempenho individual 
de jogadores. Neste trabalho, entende-se por desempenho a pontuação esperada 
de uma equipe ao longo de uma temporada de uma competição esportiva.

Entretanto, o objetivo central não se limita ao desenvolvimento 
do modelo preditivo em si. Pretende-se que tal modelo seja utilizado 
como componente de algoritmos de otimização voltados ao apoio à 
tomada de decisão na montagem de elencos. Em particular, seja $f$ o 
modelo de AM obtido e seja $X$ um arranjo contendo as \textit{features} 
correspondentes ao elenco de um determinado time. Embora o modelo possa 
ser utilizado para prever o desempenho de equipes existentes, uma vez 
treinado ele também pode ser empregado para estimar o desempenho de 
equipes hipotéticas.

Dessa forma, considere $P$ o conjunto de todos os elencos $X'$ que 
podem ser obtidos a partir do elenco atual $X$ por meio da contratação 
de novos jogadores, respeitando um determinado orçamento disponível. 
O objetivo passa então a ser identificar, dentre todas as combinações 
possíveis de jogadores em $P$, aquela que maximiza o desempenho 
previsto pelo modelo de AM. Formalmente, busca-se resolver o problema

Formalmente, busca-se resolver o problema definido pela Equação~\ref{eq:otimizacao_elenco}:

\begin{equation}
\max_{X' \in P} f(X')
\label{eq:otimizacao_elenco}
\end{equation}

Observa-se, portanto, que o problema estudado configura um problema 
de otimização combinatória, no qual o modelo de AM desempenha o papel 
de função objetivo.

\section{Objetivos Específicos}

Com o intuito de atingir o objetivo geral, foram estabelecidos os 
seguintes objetivos específicos:

\begin{enumerate}
    \item Coletar, estruturar e processar dados históricos de 
    jogadores e partidas, abrangendo múltiplas temporadas, 
    garantindo a qualidade e a integridade das informações utilizadas 
    no desenvolvimento dos modelos preditivos.
    \item Desenvolver uma metodologia para quantificar o impacto 
    individual dos jogadores, combinando estatísticas históricas ao 
    longo das temporadas por meio de um sistema de ponderação temporal, 
    priorizando desempenhos mais recentes.
    \item Construir modelos de AM para prever o desempenho agregado 
    de uma equipe em uma determinada temporada, considerando a 
    composição do elenco e os dados históricos dos jogadores.
    \item Validar a eficácia dos modelos preditivos, comparando 
    as estimativas com os resultados reais de temporadas subsequentes, 
    utilizando métricas estatísticas como erro médio absoluto (MAE) e 
    coeficiente de determinação ($R^2$).
    \item Desenvolver modelos simples de otimização baseados no 
    modelo de AM selecionado, visando a identificação de combinações 
    de jogadores que maximizem o desempenho esperado da equipe 
    sob restrições de mercado.
\end{enumerate}

Ao integrar técnicas de AM e otimização, este estudo busca 
fornecer uma ferramenta analítica de apoio à decisão para a 
montagem de elencos no futebol profissional, contribuindo para 
práticas de gestão esportiva orientadas por dados.